\begingroup
\sloppy

\section{Current object and service map}\label{modern_objects}
\index{Object design!current services}

Figure~\ref{hierarchy} includes both the inheritance hierarchy and the modern
services that compose experiments around it.  A component-shaped box is not a
new base class: it is a namespace service or helper object with one explicit
responsibility.  The arrows describe use or ownership, not inheritance unless
the label says ``base of.''

The important division is between mechanisms and coordination.  \texttt{DataSet}
owns data and materializes row partitions; \texttt{nsplit} chooses indices;
\texttt{Network} objects fit models; \texttt{obd} selects an architecture;
\texttt{crossval} repeats a supplied procedure; \texttt{evaldesign} records what
design was achieved and which inference it permits; the AUC covariance modules
estimate uncertainty; and \texttt{cvreport} renders and writes results.  No
outer coordinator reimplements a lower layer's rule.

\subsection{Current additions to DataSet}
\index{DataSet!validation set}
\index{DataSet!strataKey}
\index{DataSet!groupKey}
\index{DataSet!makeFold}
\index{DataSet!randomize3}
\index{DataSet!monitorSet}
\index{DataSet!MonitorSet}
\index{DataSet!validation set}
\index{DataSet!covariate stratification}
\index{DataSet!group columns}

The original \texttt{DataSet} contract earlier in this chapter remains valid.
The following public methods extend it for validation and planned partitions:

\begin{itemize}
\item \texttt{valLoaded()}, \texttt{getValMatrix()}, and \texttt{getNumVal()}
      expose the optional validation set.
\item \texttt{monitorSet()} returns \texttt{MONITOR\_NONE},
      \texttt{MONITOR\_VALIDATION}, or \texttt{MONITOR\_TEST};
      \texttt{monitorSetName()} returns the same choice for reports.
\item \texttt{randomize3(nTest,nVal)} and
      \texttt{randomize3D(testRatio,valRatio)} create an outcome-stratified
      train/validation/test split.
\item \texttt{makeFold(trainRows,testRows,valRows)} materializes explicit raw-row
      indices and computes every scale from the training rows only.
\item \texttt{setStrataColumns}, \texttt{getStrataColumns},
      \texttt{setStrataBins}, and \texttt{getStrataBins} configure
      outcome-by-covariate splitting.
\item \texttt{setGroupColumns} and \texttt{getGroupColumns} configure exact
      joint-value groups which may not be divided.
\item \texttt{strataKey(columns,bins)} returns dense outcome-by-covariate cell
      identifiers; \texttt{groupKey(columns)} returns dense exact-match group
      identifiers.  Both read the raw matrix without changing configuration.
\end{itemize}

\subsubsection{Example}
\begin{verbatim}
vector<unsigned> group = data.groupKey({2, 3, 4});
nsplit::GroupFoldPlan p =
    nsplit::stratifiedGroupKFold(label, group, 5);
data.makeFold(trainRows, testRows, validationRows);
cout << data.monitorSetName();       // "validation"
\end{verbatim}

\subsubsection{Notes}
Column positions passed by C++ callers are zero-based input positions.  Public
REST fields are one-based and are validated and converted by the handler.
Grouping takes precedence over covariate stratification in the ordinary raw
split.  A three-way split is outcome-stratified and does not currently compose
with either policy.

\section{Current contracts added to established objects}\label{current_established_objects}
\index{TwoSet!current ROC contract}
\index{TwoSet!invalidate}
\index{TwoSet!getTrapROCarea}
\index{TwoSet!getROCfit}
\index{TwoSet!bootstrapROC}
\index{TwoSet!confidence interval}
\index{TwoSet!operating points}
\index{TwoSet!Pearson statistic}
\index{Model!error objective query}
\index{Iterative!current training contract}
\index{Iterative!StopReason}
\index{Iterative!converged}
\index{Iterative!Observer}
\index{Iterative!quiet training}
\index{Iterative!plateau stopping}
\index{Network!current diagnostics}
\index{Network!sampleTestError}
\index{Network!condition number}
\index{SimpleProp!hidden-layer resizing}
\index{SimpleProp!growHidden}
\index{SimpleProp!removeHidden}
\index{Logistic regression!structured results}
\index{Logistic regression!coefficients}
\index{Logistic regression!Wald test}
\index{RegressNet!structured results}
\index{RegressNet!progress}
\index{RegressNet!candidate audit}
\index{RegressNet!completion status}

The detailed specifications later in this chapter describe the established
objects.  The following additions are part of their current public contracts and
must be read with those specifications.

\subsection{Model}
The following method exposes the objective needed by likelihood-based callers:

\begin{itemize}
\item \texttt{bool getXEerror() const}: returns \texttt{true} when
      cross-entropy is active and \texttt{false} for least-mean-square error.
\end{itemize}

Callers whose mathematics depends on likelihood, notably stepwise Wilks
comparisons, query this state and refuse an incompatible source fit; they do not
mutate the model's objective to make it appear compatible.

\subsection{TwoSet}
\texttt{TwoSet} now exposes explicit cache invalidation and structured ROC
results.  The following methods are available:

\begin{itemize}
\item \texttt{void invalidate()}: clears every cached statistic after a caller
      changes predictions through a mutable accessor.
\item \texttt{unsigned getTP()}, \texttt{getTN()}, \texttt{getFP()}, and
      \texttt{getFN()}: return the four confusion-matrix counts at the current
      threshold.
\item \texttt{double getTrapROCarea()}: returns the exact empirical
      Mann--Whitney AUC over every distinct operating point.
\item \texttt{const vector<double>\& getROCx()} and \texttt{getROCy()}:
      return the empirical ROC coordinates, $1-$specificity and sensitivity.
\item \texttt{ROCfit getROCfit()}: returns the current binormal ROC fit as a
      structured value; an unsupported fit has \texttt{valid=false}.
\item \texttt{void bootstrapROC()}: performs the configured number of
      outcome-stratified bootstrap resamples and caches a percentile interval
      for the binormal area.
\item \texttt{CI getStatCi()}: returns the cached bootstrap interval.
\item \texttt{void setBootstrapResamples(unsigned n)}: sets the requested
      resample count; zero disables the bootstrap.
\item \texttt{unsigned getBootstrapResamples()}: returns that count.
\item \texttt{unsigned getStatPoints()}: returns the number of distinct
      interior operating points fitted by the zROC line.
\item \texttt{double getStatChi2()}: returns the zROC fit chi-square.
\item \texttt{double getPearsonX2()}: returns the individual-level Pearson
      statistic; it is not assigned a chi-square p-value on continuous data.
\end{itemize}

\texttt{ROCfit} contains the following public variables:

\begin{itemize}
\item \texttt{double az}: fitted binormal ROC area.
\item \texttt{double chi2}: zROC fit chi-square, or not-a-number when it cannot
      be computed.
\item \texttt{double p}: fit p-value, or not-a-number when unavailable.
\item \texttt{unsigned points}: operating points included in the fit.
\item \texttt{bool valid}: whether a fit was computed.
\end{itemize}

\texttt{CI} contains:

\begin{itemize}
\item \texttt{double lo, hi}: 2.5th and 97.5th percentiles.
\item \texttt{double se}: standard deviation of successful resampled areas.
\item \texttt{unsigned resamples}: successful resamples.
\item \texttt{unsigned failures}: resamples for which no area was computable.
\item \texttt{bool valid}: whether enough resamples succeeded.
\end{itemize}

\subsection{Iterative}
\texttt{StopReason} identifies why the last run ended:

\begin{itemize}
\item \texttt{STOP\_NONE}: training has not run.
\item \texttt{STOP\_MIN\_ERROR}, \texttt{STOP\_CHANGE},
      \texttt{STOP\_WINDOW}, \texttt{STOP\_GRADMAX},
      \texttt{STOP\_PLATEAU}, and \texttt{STOP\_EARLY\_STOP}: a configured
      stopping rule fired, so the run converged under its declared rule.
\item \texttt{STOP\_MAX\_ITERATIONS}: the safety ceiling was exhausted; this is
      not convergence.
\item \texttt{STOP\_CANCELLED}: an external request stopped the fit.
\item \texttt{STOP\_PROBE\_BUDGET}: an optimizer probe used its wall-clock
      allowance; the probe is an experiment, not a completed fit.
\end{itemize}

The following methods expose and control this state:

\begin{itemize}
\item \texttt{static const char* stopReasonToken(StopReason r)}: returns the
      single machine-readable name used in reports and JSON.
\item \texttt{static bool converged(StopReason r)}: returns \texttt{true} only
      for the six stopping-rule outcomes listed above.
\item \texttt{void setObserver(Observer* obs)} and
      \texttt{Observer* getObserver()}: install or return the non-owning
      iteration observer.
\item \texttt{StopReason getStopReason()}: returns the last run's reason.
\item \texttt{unsigned getIterations()}: returns iterations completed by the
      last run.
\item \texttt{void setAutoStop(bool flag,double tol,unsigned win)}: enables or
      disables plateau stopping and sets its tolerance and window.
\item \texttt{bool getAutoStop()}, \texttt{double getAutoStopTol()}, and
      \texttt{unsigned getAutoStopWindow()}: return that configuration.
\item \texttt{void setQuiet(bool flag)} and \texttt{bool getQuiet()}:
      suppress or query narration and the reporting epilogue.  Quiet mode never
      suppresses a calculation used by fitting or stopping.
\end{itemize}

An \texttt{Iterative::Observer} supplies:

\begin{itemize}
\item \texttt{bool onIteration(unsigned iteration,double setError)}: receives
      every completed iteration and returns \texttt{false} to stop.
\item \texttt{StopReason whyStopped() const}: identifies why the observer
      returned \texttt{false}; the conservative default is cancellation.
\end{itemize}

\subsection{Network and concrete models}
The following current methods supplement the detailed \texttt{Network}
specification later in this chapter:

\begin{itemize}
\item \texttt{double sampleTestError(unsigned stride)}: returns mean error over
      every \texttt{stride}-th exemplar in the authoritative validation/test
      monitoring set without changing cached statistics; returns $-1$ when no
      one-output monitoring set exists.
\item \texttt{double getCondNum() const}: returns the last computed condition
      number, or $-1$ before calculation.
\item \texttt{double getCondMaxEig() const} and
      \texttt{getCondMinEig() const}: return the extreme eigenvalues used in
      that condition number.
\end{itemize}

\texttt{OneHiddenNet} sits between \texttt{Network} and the two
one-hidden-layer models.
\index{OneHiddenNet}
\index{OneHiddenNet!setHidden}
\index{OneHiddenNet!randomize}
\index{OneHiddenNet!save}
\index{OneHiddenNet!load}
\index{OneHiddenNet!innerTrainSet}
\index{OneHiddenNet!propagate}
\index{Forward propagation!shared implementation}
\index{SimpleProp!shared implementation}
\index{BareProp!shared implementation}
\texttt{SimpleProp} and \texttt{BareProp} describe the same architecture ---
one hidden layer and one output node --- differing only in whether bias nodes
are present.  They therefore carry identical state, and several of their
methods differed by nothing but the class name.  \texttt{OneHiddenNet} owns
that state (the hidden weight, update, and gradient matrices; the input,
hidden-output, output-weight, hidden-error, and output-gradient vectors; the
output error term; and the hidden-unit count) and implements what does not
depend on the bias architecture:

\begin{itemize}
\item \texttt{virtual void setHidden(unsigned n)}: pure virtual.  Sizing is the
      bias architecture, so each concrete model sizes its own structures;
      \texttt{load} calls this method rather than deciding widths itself.
\item \texttt{virtual void randomize()}: assigns uniform random weights within
      the configured limit and marks the weights set.
\item \texttt{virtual bool save(string\& filename)}: writes the model file ---
      the concrete class's \texttt{outputHeader}, the hidden layer, then the
      output layer.
\item \texttt{virtual bool load(string\& filename)}: reads that file back,
      refusing a file whose input count disagrees with the loaded dataset.
\item \texttt{virtual void pack()}: flattens the hidden and output gradients
      into the packed gradient vector the conjugate-direction optimizers use.
\item \texttt{virtual double innerTrainSet()}: performs one training pass over
      the training set and returns the set error.  This is the training loop
      itself --- on-line and batch, canonical descent and the separate-gradient
      path, weight decay, and the optimizer call --- for both models.
\item \texttt{void propagate()}: applies the forward equations to the exemplar
      already in the input vector, with the hidden bias slot already
      established if the model has one.  Non-virtual, and called directly by
      each concrete \texttt{forward()}: this is a per-exemplar path and may not
      acquire indirect calls in order to remove source duplication.
\end{itemize}

The training loop and the forward equations are shared with no
parameterization at all, rather than with a bias flag.  Their one differing
statement, the hidden error term, is the same formula once its domain is
written out: elements $0$ through $n_{\mathrm{hidden}}-1$ are the hidden units,
which is every element of an unbiased hidden-output vector and all but the
pinned bias slot of a biased one.  The ranged operations of \texttt{Matrix} and
the vector operators express that domain directly.

Three things deliberately remain in the concrete classes.  \texttt{df()} is
two published parameter counts, $((n_{\mathrm{input}}+2)n_{\mathrm{hidden}})+1$
with biases and $(n_{\mathrm{input}}+1)n_{\mathrm{hidden}}$ without.
\texttt{outputHeader()} writes a model-file format line that \texttt{load}
reads back.  Every width that follows from the bias column and the pinned bias
slot --- \texttt{setDataSet}, \texttt{setHidden}, \texttt{unpack},
\texttt{removeInputs}, and the statements of \texttt{forward()} that read the
exemplar and, in \texttt{SimpleProp} alone, pin the bias slot --- stays where
the biased and unbiased architectures can be read side by side.

Two notes matter to callers.  \texttt{SimpleProp} and \texttt{BareProp} remain
distinct concrete types: type identity drives cloning, the model factory, and
the first line of every saved model file, and neither is assignable to the
other.  Nothing in the shared implementation reads
\texttt{Network::biasFlag}; that flag is publicly settable, and a shared method
that consulted it could reshape a live model or change what it writes to disk.
The bias architecture is a property of the concrete type.

\texttt{SimpleProp} adds the following OBD mechanisms:

\begin{itemize}
\item \texttt{void growHidden(unsigned extra)}: adds hidden units as a warm
      start, preserving existing weights and assigning each new unit zero
      outgoing contribution.
\item \texttt{void removeHidden(const vector<unsigned>\& indices)}: removes
      the named non-bias hidden units while preserving the order and weights of
      survivors; at least one unit must remain.
\item \texttt{vector<double> hiddenSaliency()}: returns one saliency per
      non-bias hidden unit, $|w_j|$ times the training-set standard deviation of
      that unit's output.
\end{itemize}

\texttt{Logistic} exposes its last coefficient audit:

\begin{itemize}
\item \texttt{const vector<double>\& getBetas() const}: returns fitted
      coefficients; the last element is the intercept.
\item \texttt{const vector<double>\& getBetaSE() const}: returns coefficient
      standard errors; $-1$ marks an unavailable value.
\item \texttt{const vector<double>\& getWaldP() const}: returns coefficient
      Wald-test p-values.
\end{itemize}

\subsection{DFA, LDFA and QDFA}
\index{DFA}
\index{DFA!train}
\index{DFA!fitDiscriminant}
\index{LDFA!linear discriminant}
\index{QDFA!quadratic discriminant}
\index{Discriminant function analysis!shared scaffold}
\texttt{DFA} is the discriminant-analysis base.  It owns the class matrices,
mean vectors, a priori probabilities and constants that both concrete analyses
build, and since 2026 it also owns the run itself:

\begin{itemize}
\item \texttt{double train() override final}: performs one analysis.  It writes
      the banner and header, calls the fit and then the report, catches a
      singular covariance and reports the refusal, writes the history and
      last-operation files, and returns $-1$ because a discriminant analysis has
      no set error.  It is \texttt{final} because no subclass may replace the
      scaffold, and it remains a virtual override of \texttt{Model::train()} ---
      cross-validation reaches it through a \texttt{unique\_ptr<Model>}.
\item \texttt{virtual void fitDiscriminant()}: pure virtual, protected.  The
      fit, and only the fit.  Runs once per analysis, inside the scaffold's
      \texttt{try}.
\item \texttt{void reportAccuracy(ostream\&) override = 0}: pure virtual.  Each
      model writes this out in full.
\end{itemize}

\texttt{LDFA} fits one \emph{pooled} covariance across all classes, inverts it,
and forms constants $K_c = \frac{1}{2}\mathbf{u}_c^{\mathsf{T}} S
\mathbf{u}_c - \log P_c$.  Its discriminant is $d_c = \mathbf{u}_c^{\mathsf{T}}
S \mathbf{x} - K_c$ and the \emph{larger} discriminant wins; the binary graded
score is $\sigma(d_1 - d_0)$.

\texttt{QDFA} fits a \emph{separate} covariance per class, inverts each, and
forms constants $K_c = \log|C_c| - \log P_c$.  Its discriminant is $d_c =
(\mathbf{x}-\mathbf{u}_c)^{\mathsf{T}} S_c (\mathbf{x}-\mathbf{u}_c) + K_c$ and
the \emph{smaller} discriminant wins, because it is a distance where the linear
one is a similarity; the binary graded score is $\sigma(d_0 - d_1)$.

Those opposite senses are why only the scaffold is shared.  Sharing the
per-exemplar scoring loops would require either a virtual call per exemplar or a
comparator parameter, and the base class deliberately contains no sign, winner
flag, formula descriptor or type switch: the published mathematics stays in the
class that owns it.  Both models write a \emph{graded} score through the
logistic function rather than a hard decision, so the ROC sweep sees a curve;
the classification boundary is unchanged.

\subsection{RegressNet}
The following methods support non-interactive, cancellable stepwise analysis:

\begin{itemize}
\item \texttt{void setThreshold(double t)}: supplies the p-value threshold
      without an interactive prompt.
\item \texttt{void setProgress(ProgressFn f)}: installs the progress callback.
\item \texttt{void setObserver(Iterative::Observer* obs)}: installs the
      non-owning cancellation observer on every candidate clone.
\item \texttt{const vector<Candidate>\& getCandidates() const}: returns every
      candidate in evaluation order.
\item \texttt{const vector<pair<double,unsigned> >\& getSelectionPath() const}:
      returns selected p-value/variable pairs in selection order.
\item \texttt{const vector<unsigned>\& getFinalVariables() const}: returns the
      variables settled by completed passes.
\item \texttt{unsigned getFitsCompleted() const}: returns candidate fits made.
\item \texttt{bool getComplete() const}: distinguishes a completed selection
      from a partial audit after failure or cancellation.
\end{itemize}

\texttt{Progress} contains direction, pass number, candidate number and pass
total, completed-fit count, conceptual variable, its full input-node group,
phase, completion flag, convergence flag, and stop-reason token.

\texttt{Candidate} contains pass/candidate/variable identity, the full input
group, prior and candidate errors, degrees of freedom, $G^2$, p-value,
convergence, stop reason, iterations, and whether it won its pass.

\subsubsection{Example}
\begin{verbatim}
network.setAutoStop(true, 1e-4, 100);
double error = network.train();
if (!Iterative::converged(network.getStopReason())) rejectFit();

TwoSet::ROCfit fit = data.getTestTwoSet().getROCfit();
const vector<double>& p = logistic.getWaldP();
\end{verbatim}

\subsubsection{Notes}
Reaching \texttt{STOP\_MAX\_ITERATIONS} is operation completion but not fit
convergence.  Quiet mode changes output only.  The ROC bootstrap uses an
independent resampling stream and cannot perturb training.  Stepwise candidate
fits must converge under cross-entropy before entering a Wilks comparison, and
the original model is never mutated.

\section{Partition planning: nsplit}\label{nsplit_objects}
\index{nsplit}
\index{Group-disjoint partition}
\index{Stratified partition}

The \texttt{nsplit} namespace is an index-level service.  It never scales data,
fits a model, or interprets an input column.

\texttt{Holdout} contains:

\begin{itemize}
\item \texttt{vector<unsigned> test, train}: row positions assigned to each
      side.
\item \texttt{unsigned n0, n1}: whole-sample outcome counts.
\item \texttt{unsigned n0Test, n1Test}: achieved test-set outcome counts.
\end{itemize}

\texttt{StratHoldout} contains:

\begin{itemize}
\item \texttt{vector<unsigned> test, train}: row positions assigned to each
      side.
\item \texttt{vector<unsigned> cellTotal}: whole-sample count in each stratum.
\item \texttt{vector<unsigned> cellTest}: achieved test count in each stratum.
\end{itemize}

\texttt{GroupHoldout} contains:

\begin{itemize}
\item \texttt{vector<unsigned> test, train}: row positions assigned to each
      side, with every group wholly on one side.
\item \texttt{unsigned nGroups}: number of distinct groups.
\item \texttt{vector<unsigned> groupsInTest}: group identifiers assigned to
      the test side.
\end{itemize}

\texttt{FoldPlan} contains:

\begin{itemize}
\item \texttt{bool ok} and \texttt{string reason}: planning status and refusal.
\item \texttt{vector<unsigned> foldId}: fold number for every input row.
\item \texttt{unsigned k, nStrata}: requested folds and observed strata.
\item \texttt{vector<unsigned> foldSize}: achieved rows per fold.
\item \texttt{vector<vector<unsigned> > stratumByFold}: achieved stratum counts
      by fold.
\item \texttt{vector<string> warnings}: non-fatal planning facts.
\end{itemize}

\texttt{GroupFoldPlan} extends \texttt{FoldPlan} with:

\begin{itemize}
\item \texttt{unsigned nGroups}: distinct groups in the planned rows.
\item \texttt{vector<unsigned> groupsPerFold}: whole groups assigned per fold.
\item \texttt{vector<array<unsigned,2> > classByFold}: negative and positive
      rows achieved per fold.
\item \texttt{unsigned leakageCount}: groups found in more than one fold; a
      valid plan has zero.
\item \texttt{double imbalanceScore}: worst relative row or class imbalance.
\item \texttt{unsigned largestGroup}: size of the largest indivisible group.
\end{itemize}

The following functions are available:

\begin{itemize}
\item \texttt{string partitionError(unsigned nRows,train,test,bool coverage)}:
      returns an empty string for a valid two-way row partition, otherwise a
      reason describing an out-of-range, duplicate, overlapping, or missing row.
\item \texttt{Holdout stratifiedHoldout(label,nTest)}: returns an
      outcome-stratified train/test partition of binary labels.
\item \texttt{StratHoldout holdoutByStrata(stratum,nTest)}: apportions test
      positions across arbitrary stratum identifiers by largest remainder.
\item \texttt{GroupHoldout groupHoldout(label,group,nTest)}: assigns each
      complete group to training or test while approaching the requested size
      and outcome balance.
\item \texttt{FoldPlan stratifiedKFold(stratum,k)}: returns a balanced k-fold
      assignment over arbitrary stratum identifiers.
\item \texttt{GroupFoldPlan stratifiedGroupKFold(label,group,k)}: returns an
      outcome-balanced whole-group fold assignment with measured leakage and
      imbalance.
\item \texttt{vector<unsigned> kFold(label,k)}: returns the outcome-only case
      of \texttt{stratifiedKFold}.
\end{itemize}
\index{nsplit!stratifiedGroupKFold}
\index{nsplit!Holdout}
\index{nsplit!StratHoldout}
\index{nsplit!GroupHoldout}
\index{nsplit!FoldPlan}
\index{nsplit!GroupFoldPlan}
\index{nsplit!partitionError}
\index{nsplit!stratifiedHoldout}
\index{nsplit!holdoutByStrata}
\index{nsplit!groupHoldout}
\index{nsplit!stratifiedKFold}
\index{nsplit!kFold}
\index{FoldPlan}
\index{GroupFoldPlan}
\index{Leakage!partition validation}

\subsubsection{Example}
\begin{verbatim}
util::set_seed(42);
nsplit::GroupFoldPlan plan =
    nsplit::stratifiedGroupKFold(label, group, 5);
if (!plan.ok || plan.leakageCount != 0)
    throw runtime_error(plan.reason);
\end{verbatim}

\subsubsection{Notes}
Every plan is reproducible under the engine seed.  Group identifiers are merely
equivalence labels; renumbering them must not change the partition.  Achieved
fold counts and leakage are recomputed from the final assignment.  Indivisible
groups may make exact sizes or exact class balance impossible, so imbalance is
reported rather than hidden.  Fewer than \texttt{k} groups is a refusal.

\section{Typed evaluation design: evaldesign}\label{evaldesign_objects}
\index{evaldesign}
\index{Sampling unit}
\index{AUC inference policy}
\index{evaldesign!Partition}
\index{evaldesign!InferenceChoice}
\index{evaldesign!SamplingUnit}
\index{evaldesign!PartitionMethod}
\index{evaldesign!chooseInference}
\index{Evaluation design!typed vocabulary}

\texttt{evaldesign} is the vocabulary and policy layer; it computes no
statistic.  Its enumerations are:

\begin{itemize}
\item \texttt{SamplingUnit}: \texttt{Unspecified}, \texttt{Row}, or
      \texttt{Cluster}, declaring the independent sampling unit.
\item \texttt{AucInference}: \texttt{None}, \texttt{DeLongIndependent}, or
      \texttt{ObuchowskiClustered}, naming the permitted estimator.
\item \texttt{PartitionMethod}: \texttt{OutcomeStratified},
      \texttt{CovariateStratified}, or \texttt{StratifiedGroup}, naming how
      rows were assigned.
\end{itemize}

\texttt{Partition} contains the following public variables:

\begin{itemize}
\item \texttt{PartitionMethod method}: assignment method.
\item \texttt{unsigned seed}: reproducibility seed.
\item \texttt{unsigned k}: number of folds; zero for a holdout.
\item \texttt{vector<unsigned> strataColumns}: caller's one-based covariate
      columns; \texttt{strataBins} and \texttt{nStrata} give bin count and
      achieved cell count.
\item \texttt{vector<unsigned> groupColumns}: caller's one-based group columns;
      \texttt{nGroups} gives distinct groups in the partitioned rows.
\item \texttt{bool developmentOnly}: whether the folds exclude a locked test.
\item \texttt{unsigned nRequested, nAchieved}: requested and achieved holdout
      sizes; indivisible groups may make them differ.
\item \texttt{unsigned leakage}: groups or rows found on both sides; a valid
      group plan requires zero.
\item \texttt{vector<string> warnings}: non-fatal design facts.
\item \texttt{string refusal}: reason no partition could be built.
\item \texttt{foldRows, foldEvents, foldGroups}: parallel achieved counts for
      each fold.
\item \texttt{double imbalanceScore}: worst relative deviation from a fold's
      fair row or class share.
\item \texttt{unsigned largestGroup}: largest indivisible cluster size.
\end{itemize}

\texttt{InferenceChoice} contains \texttt{AucInference method} and
\texttt{string reason}; the reason is required when the method is
\texttt{None}.

The following functions are available:

\begin{itemize}
\item \texttt{bool parseSamplingUnit(string text,SamplingUnit\& out)}: parses
      the REST value \texttt{rows}, \texttt{cluster}, or an empty declaration;
      returns \texttt{false} for an unrecognized token.
\item \texttt{samplingUnitName}, \texttt{samplingUnitPhrase}, and
      \texttt{independenceStatus}: return the authoritative display wording for
      a sampling-unit declaration.
\item \texttt{inferenceName}, \texttt{inferenceText}, and
      \texttt{inferenceShortName}: return long, reason-bearing, and compact
      names for an inference choice.
\item \texttt{partitionMethodName}: returns the display name of a partition
      method.
\item \texttt{string describeFolds(const Partition\& p)}: derives a fold-plan
      sentence from structured fields.
\item \texttt{string describeHoldout(const Partition\& p)}: derives the
      corresponding holdout sentence, including requested and achieved sizes.
\item \texttt{InferenceChoice chooseInference(unit,split,estimable,reason)}:
      returns the only covariance estimator permitted by the declared unit and
      achieved split, or \texttt{None} with a reason.
\end{itemize}

Display strings must come from these functions rather than being composed by a
GUI or report.

\subsubsection{Example}
\begin{verbatim}
evaldesign::InferenceChoice c = evaldesign::chooseInference(
    evaldesign::SamplingUnit::Cluster,
    evaldesign::PartitionMethod::StratifiedGroup,
    clustered.ok, clustered.reason);
\end{verbatim}

\subsubsection{Notes}
A group-disjoint split prevents leakage; it does not establish row independence.
An unspecified unit intentionally withholds inferential claims.  Independent-row
inference over a grouped design is refused, and neither covariance estimator may
silently substitute for the other.

\section{Shared AUC covariance algebra: auccov}\label{auccov_objects}
\index{auccov}
\index{ROC area!placements}
\index{Mann--Whitney area}
\index{auccov!midranks}
\index{auccov!Placements}
\index{auccov!Interval}
\index{auccov!Contrast}
\index{auccov!compute}
\index{auccov!interval}
\index{auccov!contrast}
\index{ROC area!ties}
\index{ROC area!zero-variance contrast}

\texttt{auccov} owns the operations that are invariant to the sampling unit.
\texttt{Placements} contains:

\begin{itemize}
\item \texttt{bool ok} and \texttt{string reason}: computation status.
\item \texttt{unsigned n0, n1}: negative and positive row counts.
\item \texttt{vector<unsigned> pos, neg}: original row indices by outcome.
\item \texttt{vector<vector<double> > v10, v01}: procedure-by-observation
      structural components.
\item \texttt{vector<double> auc}: exact point AUC for each procedure.
\end{itemize}

\texttt{Interval} contains point \texttt{auc}, standard error \texttt{se},
limits \texttt{lo}/\texttt{hi}, and \texttt{valid}.

\texttt{Contrast} contains the two areas and their standard errors, difference
\texttt{delta} and \texttt{seDelta}, each area's interval, the difference
interval, \texttt{z}, two-sided \texttt{p}, zero-variance flags
\texttt{degenerate}/\texttt{separated}, explanatory \texttt{note}, and
\texttt{valid}.

The following functions are available:

\begin{itemize}
\item \texttt{vector<double> midranks(const vector<double>\& x)}: returns
      one-based ranks; tied observations receive the mean rank they span.
\item \texttt{Placements compute(label,pred)}: returns positive/negative row
      indices, structural components, and exact point AUCs for paired procedures.
\item \texttt{Interval interval(auc,cov,i,bool clamp)}: returns the normal
      interval for area \texttt{i}; \texttt{clamp} controls truncation to
      $[0,1]$.
\item \texttt{Contrast contrast(auc,cov,i,j,bool clamp)}: returns
      $AUC_i-AUC_j$, its standard error, interval, z statistic, and two-sided
      p-value.
\end{itemize}

The estimator above this layer supplies the covariance matrix; it does not
reimplement these operations.

\subsubsection{Example}
\begin{verbatim}
auccov::Placements p = auccov::compute(label, predictions);
auccov::Contrast c =
    auccov::contrast(p.auc, covariance, 0, 1, false);
\end{verbatim}

\subsubsection{Notes}
Ties receive half credit through common mid-ranks, so the ordinary and clustered
paths have identical point AUCs.  A zero-variance, zero-difference contrast has
\texttt{p=1}; a zero-variance nonzero difference is deterministic separation
with \texttt{p=0}.  A materially negative difference variance is refused.

\section{AUC covariance estimators}\label{auc_estimators}
\index{DeLong covariance}
\index{delong}
\index{delong!analyze}
\index{delong!Result}
\index{delong!interval}
\index{delong!contrast}
\index{Obuchowski covariance}
\index{clustered auc@\texttt{clustered\_auc}}
\index{clustered auc@\texttt{clustered\_auc}!analyze}
\index{clustered auc@\texttt{clustered\_auc}!Result}
\index{clustered auc@\texttt{clustered\_auc}!interval}
\index{clustered auc@\texttt{clustered\_auc}!contrast}
\index{Clustered ROC data}

\subsection{delong}
\texttt{delong::Result} contains:

\begin{itemize}
\item \texttt{unsigned n0, n1}: negative and positive row counts.
\item \texttt{vector<double> auc}: one exact point area per procedure.
\item \texttt{Matrix<double> cov}: paired DeLong covariance matrix.
\item \texttt{bool ok} and \texttt{string reason}: estimator status and refusal.
\end{itemize}

The following functions are available:

\begin{itemize}
\item \texttt{Result analyze(label,pred)}: estimates paired DeLong covariance
      for procedures scored on the same independent rows.
\item \texttt{Interval interval(const Result\& r,unsigned i)}: returns the
      normal DeLong interval for procedure \texttt{i}, clamped to $[0,1]$.
\item \texttt{Contrast contrast(const Result\& r,unsigned i,unsigned j)}:
      returns the paired area contrast $AUC_i-AUC_j$.
\end{itemize}

\subsection{clustered\_auc}
\texttt{clustered\_auc::Result} contains:

\begin{itemize}
\item \texttt{bool ok} and \texttt{string reason}: estimator status and refusal.
\item \texttt{unsigned nRows, nClusters}: row and independent-cluster counts.
\item \texttt{unsigned clusters10, clusters01}: clusters carrying at least one
      positive or negative row; these are the covariance divisors' populations.
\item \texttt{unsigned n0, n1}: negative and positive row counts.
\item \texttt{vector<double> auc}: exact patient-row point areas.
\item \texttt{Matrix<double> cov}: Obuchowski clustered covariance matrix.
\item \texttt{string method}: authoritative estimator name for reports.
\end{itemize}

The following functions are available:

\begin{itemize}
\item \texttt{Result analyze(label,cluster,pred)}: estimates Obuchowski's
      nonparametric clustered covariance for paired procedures.
\item \texttt{Interval interval(const Result\& r,unsigned i)}: returns the
      unclamped published normal interval for procedure \texttt{i}.
\item \texttt{Contrast contrast(const Result\& r,unsigned i,unsigned j)}:
      returns the clustered paired area contrast.
\end{itemize}

\subsubsection{Example}
\begin{verbatim}
delong::Result iid = delong::analyze(label, predictions);
clustered_auc::Result grouped =
    clustered_auc::analyze(label, cluster, predictions);
// The design policy decides which result may be reported.
\end{verbatim}

\subsubsection{Notes}
DeLong requires at least two rows of each outcome class and assumes independent
rows.  Obuchowski requires at least two clusters carrying each outcome class.
Clustering can increase or decrease the estimated standard error; its defining
property is that the cluster is the independent unit.  Out-of-fold predictions
are not an independent test sample for either locked-test contrast.

\section{Training-control services}\label{training_control_objects}

\subsection{PlateauDetector}
\index{PlateauDetector}
\index{Early stopping!plateau}
\index{PlateauDetector!update}
\index{PlateauDetector!reset}
\texttt{PlateauDetector} maintains two adjacent moving-average windows.  The
following methods are available:

\begin{itemize}
\item \texttt{PlateauDetector(unsigned window,double tol,unsigned patience)}:
      constructs a detector; window is at least two and patience at least one.
\item \texttt{void reset()}: clears both windows and the strike count.
\item \texttt{bool update(double error)}: adds one iteration error and returns
      \texttt{true} after the configured number of consecutive comparisons fail
      to achieve the required relative improvement.
\item \texttt{unsigned strikes() const}: returns the current consecutive-strike
      count.
\end{itemize}

\subsection{autoalgo}
\index{autoalgo}
\index{Optimizer!automatic selection}
\index{autoalgo!Probe}
\index{autoalgo!Result}
\index{autoalgo!pick}
\index{Optimizer!probe budget}
\texttt{autoalgo::Probe} contains:

\begin{itemize}
\item \texttt{int algorithm} and \texttt{string name}: training-menu/API
      number (1 through 3) and display name.
\item \texttt{double error}: final probe error.
\item \texttt{unsigned iterations}: iterations completed within the probe.
\item \texttt{bool usable}: whether the result may compete for selection.
\item \texttt{Iterative::StopReason stop}: reason the probe ended.
\end{itemize}

\texttt{autoalgo::Result} contains:

\begin{itemize}
\item \texttt{int selected} and \texttt{string selectedName}: winning
      algorithm and display name.
\item \texttt{vector<Probe> probes}: complete optimizer audit.
\item \texttt{unique\_ptr<Network> winner}: selected clone, including probe
      progress.
\item \texttt{bool cancelled}: whether external cancellation ended selection.
\end{itemize}

\begin{itemize}
\item \texttt{Result pick(const Network\& start,unsigned budgetMs,
      const atomic<bool>* cancel)}: probes canonical, CGD, and Shanno on clones
      of identical starting weights and returns the lowest-error usable clone.
\end{itemize}

\subsection{modelfactory and netclone}
\index{modelfactory}
\index{cloneNetwork}
\index{modelfactory!build}
\index{modelfactory!Spec}
\index{modelfactory!createByTypeName}
\index{Network!cloning}
\texttt{modelfactory::Spec} contains \texttt{logistic}, \texttt{bias}, hidden
layer sizes, \texttt{xentropy}, \texttt{logLastop}, and \texttt{logHistory}.

\begin{itemize}
\item \texttt{unique\_ptr<Model> build(const Spec\& spec,DataSet\& data,
      string\& err)}: constructs and binds a fresh untrained model in the
      required order; returns null and fills \texttt{err} for an invalid request.
\item \texttt{unique\_ptr<Model> createByTypeName(string typeName,bool bias)}:
      creates the appropriate empty concrete type for a saved-network load;
      returns null for an unknown type name.
\item \texttt{unique\_ptr<Network> cloneNetwork(const Network\& source)}:
      returns a deep copy of a known concrete network, including its dataset and
      training state; returns null for an unknown concrete type.
\end{itemize}

\subsubsection{Example}
\begin{verbatim}
PlateauDetector stop(100, 1e-4, 3);
if (stop.update(trainingError)) requestStop();

modelfactory::Spec spec;
spec.logistic = true;
unique_ptr<Model> model = modelfactory::build(spec, data, error);
\end{verbatim}

\subsubsection{Notes}
A plateau is an eligible stopping rule; an iteration ceiling is not.  Automatic
optimizer probes are bounded experiments, not completed fits.  Factory
construction and cloning centralize type knowledge; training remains virtual
behavior of the concrete model.

\section{OBD architecture selection}\label{obd_objects}
\index{Optimal Brain Damage}
\index{obd}
\index{obd!Config}
\index{obd!SizeTrial}
\index{obd!Eligibility}
\index{obd!Result}
\index{obd!run}
\index{Iteration ceiling!OBD refusal}

\texttt{obd::Eligibility} distinguishes \texttt{ELIGIBLE},
\texttt{INCOMPLETE\_CEILING}, \texttt{INCOMPLETE\_CANCELLED}, and
\texttt{NUMERICAL\_FAILURE}.

\texttt{SizeTrial} contains hidden-unit count, training and held-out errors and
classification accuracies at the validation minimum, stop reason, growth/prune
phase, eligibility, and iterations.

\texttt{Config} contains:

\begin{itemize}
\item \texttt{hStart, hMax}: first and largest hidden-unit counts.
\item \texttt{iterBudget, sampleEvery}: per-size ceiling and validation sampling
      cadence.
\item \texttt{earlyStopTol, earlyStopPatience}: tolerated held-out reversal and
      consecutive samples required to stop a size.
\item \texttt{growPatience, pruneTol}: growth and pruning decisions.
\item \texttt{algorithm}: 0/1/2 for a fixed optimizer or $-1$ for automatic
      selection.
\item \texttt{plateauTol, plateauWindow}: training-error plateau backstop.
\end{itemize}

\texttt{Result} contains success/refusal message, selected hidden size, complete
trial history, owning winner, cancellation and ceiling-exhaustion flags,
resolved optimizer, and whether it was selected automatically.

The following functions are available:

\begin{itemize}
\item \texttt{Eligibility classify(StopReason stop,double score,double error)}:
      returns the authoritative trial-admission classification.
\item \texttt{const char* stopToken(StopReason stop,Eligibility e)}: returns the
      machine-readable trial outcome.
\item \texttt{Result run(DataSet\& data,const Config\& cfg,ProgressFn progress,
      const atomic<bool>* cancel)}: performs the grow-then-prune search and
      returns the owning selected network only after a valid search.
\end{itemize}

\subsubsection{Example}
\begin{verbatim}
obd::Config cfg;
cfg.hMax = 8;
cfg.algorithm = -1;                  // automatic optimizer probe
obd::Result r = obd::run(data, cfg, nullptr, nullptr);
if (r.ok) current = std::move(r.winner);
\end{verbatim}

\subsubsection{Notes}
OBD requires one discrete output and a held-out monitoring set.  Validation is
preferred; the test set is only the two-way fallback.  A needed trial that hits
its ceiling is refused and never compared.  In nested cross-validation the
whole search, including optimizer choice, is repeated inside each fold.

\section{Cross-validation objects}\label{crossval_objects}
\index{Cross-validation!objects}
\index{crossval}
\index{crossval!Metrics}
\index{crossval!FoldResult}
\index{crossval!RunResult}
\index{crossval!ProcResult}
\index{crossval!Progress}
\index{crossval!FoldSelection}
\index{crossval!ProcedureSpec}
\index{crossval!Comparison}
\index{crossval!LockedEntry}
\index{crossval!LockedResult}
\index{crossval!run}
\index{crossval!compare}
\index{crossval!evaluateOnce}
\index{Out-of-fold prediction}
\index{Deterministic substream}

\texttt{crossval} owns repetition over an immutable fold assignment.
\texttt{Metrics} contains:

\begin{itemize}
\item \texttt{unsigned n}: rows contributing to the calculation.
\item \texttt{double az}: fitted binormal AUC.
\item \texttt{double trap}: exact empirical AUC.
\item \texttt{double sens, spec, ca}: sensitivity, specificity, and
      classification accuracy at the selected threshold.
\end{itemize}

\texttt{FoldResult} contains:

\begin{itemize}
\item \texttt{unsigned fold, n}: fold number and held-out row count.
\item \texttt{bool ok} and \texttt{string reason}: fit status and refusal.
\item \texttt{Metrics metrics}: the fold's held-out statistics.
\end{itemize}

\texttt{RunResult} contains:

\begin{itemize}
\item \texttt{bool ok, cancelled} and \texttt{string message}: run status.
\item \texttt{vector<FoldResult> folds}: one audit record per fold.
\item \texttt{vector<double> prediction, outcome}: row-aligned out-of-fold
      predictions and outcomes.
\item \texttt{unsigned validFolds, pooledN}: successful folds and pooled rows.
\item \texttt{double pooledAz, pooledTrap}: pooled OOF binormal and exact AUCs.
\end{itemize}

The remaining coordination structures are:

\begin{itemize}
\item \texttt{ProcResult}: success, cancellation, refusal reason, and held-out
      predictions returned by one procedure invocation.
\item \texttt{Progress}: stage, current procedure and total procedures, current
      fold and total folds.
\item \texttt{FoldSelection}: selected hidden size and optimizer, plus whether
      automatic optimizer selection was used.
\item \texttt{ProcedureSpec}: display name, procedure callback, and optional
      selection-audit sink.
\item \texttt{Comparison}: shared outcome and fold arrays, fold count, and one
      timed \texttt{RunResult} entry per procedure.
\item \texttt{LockedEntry}: procedure name, fit status and reason, paired test
      predictions, elapsed seconds, and selection audit.
\item \texttt{LockedResult}: shared original test-row identities and outcomes,
      cancellation/message, and all \texttt{LockedEntry} values.
\end{itemize}

\texttt{Procedure} is the model-family callback which fits one materialized
training fold and returns held-out predictions.  \texttt{ProgressFn} reports the
current procedure, fold, and stage.

The following functions are available:

\begin{itemize}
\item \texttt{Metrics metricsFor(outcome,pred,rows)}: calculates held-out
      classification and ROC metrics for the named rows; unavailable metrics
      are $-1$.
\item \texttt{RunResult run(DataSet\& data,foldId,Procedure proc,...)}: runs one
      procedure over every fold and scatters successful predictions back to
      their original row positions.
\item \texttt{Comparison compare(DataSet\& data,foldId,procs,...)}: runs several
      named procedures over one immutable fold plan and returns paired OOF
      results.
\item \texttt{LockedResult evaluateOnce(DataSet\& data,trainRows,testRows,
      procs,...)}: refits each procedure once on the development rows and scores
      the common untouched locked test.
\end{itemize}

\subsection{cvadapters}
\index{cvadapters}
\index{cvadapters!trainProcedure}
\index{cvadapters!InnerSplit}
\index{cvadapters!dfaProcedure}
\index{cvadapters!nestedObdProcedure}
\index{cvadapters!innerValidationSplit}
\index{Nested cross-validation}
\texttt{InnerSplit} contains \texttt{ok}, refusal \texttt{reason}, and raw-row
\texttt{train} and \texttt{validation} vectors.  The following functions adapt
model families to the generic procedure contract:

\begin{itemize}
\item \texttt{Procedure trainProcedure(const Network\& templateNet,
      unsigned maxIter)}: clones and trains a configured network from fresh
      weights in each fold.
\item \texttt{Procedure dfaProcedure(bool quadratic)}: returns an LDFA adapter
      when false and QDFA adapter when true.
\item \texttt{Procedure nestedObdProcedure(const obd::Config\& cfg,
      double innerValFraction,...)}: performs the complete OBD architecture
      search within each outer training fold and returns only outer held-out
      predictions.
\item \texttt{InnerSplit innerValidationSplit(trainRows,rowLabel,rowGroup,
      double fraction)}: returns a stratified inner split, group-disjoint when
      \texttt{rowGroup} is supplied; returns \texttt{ok=false} rather than
      falling back when the grouped split is infeasible.
\end{itemize}

\subsubsection{Example}
\begin{verbatim}
crossval::ProcedureSpec logistic{
    "Logistic",
    cvadapters::trainProcedure(templateNet, 20000), nullptr};
crossval::Comparison c = crossval::compare(
    data, foldPlan.foldId, {logistic}, nullptr, true, 42);
\end{verbatim}

\subsubsection{Notes}
A failed fold contributes no fabricated prediction and remains visibly missing.
All procedures share the same fold plan.  Deterministic substreams key fits by
seed, procedure name, and fold, so comparison membership and order cannot change
a procedure's fit.  Nested OBD's inner validation is group-disjoint under a
grouped design; infeasibility is a failure, never a row-wise fallback.

\section{Cross-validation reporting: cvreport}\label{cvreport_objects}
\index{cvreport}
\index{Cross-validation!artifacts}
\index{cvreport!writeArtifacts}
\index{cvreport!tier1}
\index{cvreport!tier2}
\index{cvreport!render}
\index{cvreport!PlanInfo}
\index{cvreport!LockedInfo}
\index{cvreport!LockedColumn}
\index{cvreport!LockedContrast}
\index{cvreport!ArtifactResult}
\index{Raw-row identity}

\texttt{PlanInfo} contains the dataset label and whole-dataset totals, typed fold
design, primary/reference names, optional per-row cluster identity, original
raw-row identity, and original row count.  Its methods are:

\begin{itemize}
\item \texttt{string foldPlanText() const}: derives the fold-plan heading from
      the typed partition.
\item \texttt{string rawRowError(unsigned n) const}: returns an empty string for
      a valid absent/identity map or for exactly \texttt{n} distinct in-range
      raw-row identifiers; otherwise returns the defect.
\end{itemize}

\texttt{LockedColumn} contains:

\begin{itemize}
\item \texttt{string name}: procedure display name.
\item \texttt{bool hasAuc} and \texttt{double auc}: point-area availability and
      value.
\item \texttt{bool hasInterval} and \texttt{double lo, hi}: interval
      availability and limits.
\item \texttt{string architecture}: selected architecture description.
\item \texttt{string note}: refusal or qualification shown with the result.
\item \texttt{vector<double> prediction}: predictions paired to locked rows.
\end{itemize}

\texttt{LockedContrast} contains:

\begin{itemize}
\item \texttt{bool hasPoint, hasInference}: availability of the AUC difference
      and of its inferential comparison.
\item \texttt{string primary, reference}: contrasted procedure names.
\item \texttt{double delta, p}: AUC difference and two-sided p-value.
\item \texttt{bool significant}: whether the declared comparison crosses its
      significance rule.
\item \texttt{bool degenerate, separated}: zero-variance case indicators.
\item \texttt{string note}: refusal or qualification.
\end{itemize}

\texttt{LockedInfo} contains:

\begin{itemize}
\item \texttt{evaldesign::Partition split}: achieved locked-holdout design.
\item \texttt{SamplingUnit samplingUnit}: declared independent unit.
\item \texttt{AucInference inference}: covariance estimator permitted by the
      design.
\item \texttt{string inferenceReason}: explanation when inference is withheld.
\item \texttt{vector<unsigned> rawRow}: original locked-row identities.
\item \texttt{vector<double> outcome}: locked outcomes.
\item \texttt{vector<unsigned> cluster}: cluster identity for each locked row.
\item \texttt{unsigned nClusters}: distinct clusters present in the locked
      sample.
\item \texttt{vector<LockedColumn> columns}: procedure results.
\item \texttt{LockedContrast contrast}: prespecified paired comparison.
\end{itemize}

\texttt{LockedInfo} derives every displayed design string from typed state.
Its \texttt{clusterError()} method validates cluster identity before clustered
inference or artifact writing.

\texttt{ArtifactResult} contains attempted artifact name and path, success, and
error text.  The following functions are available:

\begin{itemize}
\item \texttt{string tier1(cmp,info,locked)}: returns the one-screen headline
      table and verdict.
\item \texttt{string tier2(cmp,info,locked)}: returns per-fold, selection,
      timing, design, and locked-test detail.
\item \texttt{vector<ArtifactResult> writeArtifacts(cmp,info,dir,locked)}:
      writes \texttt{cv\_predictions.csv}, \texttt{cv\_metrics.csv},
      \texttt{cv\_run.json}, and, when applicable,
      \texttt{cv\_locked\_predictions.csv}; returns one status per attempted
      file.
\item \texttt{void render(ostream\& out,cmp,info,dir,locked)}: writes Tier 2,
      then Tier 1, then the machine-readable artifacts.
\end{itemize}

\subsubsection{Example}
\begin{verbatim}
cvreport::PlanInfo info;
info.folds = achievedFoldDesign;
cvreport::LockedInfo locked;
locked.split = achievedLockedDesign;
cvreport::render(cout, comparison, info, runDirectory, locked);
\end{verbatim}

\subsubsection{Notes}
The reporter chooses neither partitions nor covariance estimators.  Every design
sentence is derived from typed state.  Original raw-row and cluster identifiers
are validated before artifacts are written so development and locked prediction
files remain joinable.  Tier 1 is the conclusion, Tier 2 is descriptive detail,
and Tier 3 is the machine-readable audit substrate.

\endgroup
